\section{Introduction}

\subsection{Motivation}

The four attention projection matrices $W_Q, W_K, W_V, W_O \in \mathbb{R}^{d \times d}$ 
together contain $4 d^2$ parameters per layer -- for $d = 768$, that's **2.36M 
parameters per layer**, dominating the attention block. Across 12 layers of 
RoBERTa, this is \textbf{28.3M parameters} just in projections.

Low-rank factorization reduces this cost. For a matrix $W \in \mathbb{R}^{d \times d}$, 
SVD replaces it with $U_r V_r$ where $U_r \in \mathbb{R}^{d \times r}$, 
$V_r \in \mathbb{R}^{r \times d}$, and $r \ll d$. Parameter count drops from 
$d^2$ to $2 r d$.

\textbf{Standard SVD compression} treats each projection independently:

\[
W_Q \to U_{Q,r} V_{Q,r}, \quad W_K \to U_{K,r} V_{K,r}, \quad W_V \to U_{V,r} V_{V,r}, \quad W_O \to U_{O,r} V_{O,r}
\]

This costs $4 \cdot 2 r d = 8 r d$ parameters.

\textbf{Key observation.} In pretrained transformers, the Q, K, V projections of a 
head are \textbf{highly correlated} -- they share a common input basis. This suggests 
that a \textbf{shared} input-space basis $V_{QKV}$ may capture the structure of all 
three with minimal loss.

\subsection{Problem Statement}

Given a pretrained transformer $\mathcal{M}$ with hidden dimension $d$ and 
head dimension $d_h$, we seek compressed projections 
$\{W_Q, W_K, W_V, W_O\}$ such that:

\[
\text{Cost}(\mathcal{A}') \ll 4 d^2, \qquad
\mathcal{A}' \approx \mathcal{A}
\]

subject to:

\begin{enumerate}
\item \textbf{Parameter reduction:} Compressed projection count $< 4 d^2$.
\item \textbf{Accuracy preservation:} Near-baseline on SST-2.
\item \textbf{Multi-architecture:} Support flat, nested, causal, and GQA models.
\item \textbf{Recovery capacity:} LoRA adapters for fine-tuning.
\item \textbf{Correctness:} Handle nested encoders without double projection.

\end{enumerate}
\subsection{Our Approach: Joint QKV SVD}

We stack the Q, K, V projections:

\[
W_{QKV} = \begin{pmatrix} W_Q \\ W_K \\ W_V \end{pmatrix} \in \mathbb{R}^{(d + 2 d_{kv}) \times d}
\]

where $d_{kv} = H_{kv} \cdot d_h$ for grouped-query attention.

We then compute a single SVD:

\[
W_{QKV} = U_{QKV} \Sigma_{QKV} V_{QKV}^\top
\]

and truncate to rank $r$:

\[
W_{QKV} \approx U_{QKV, r} \Sigma_{QKV, r}^{1/2} \cdot \Sigma_{QKV, r}^{1/2} V_{QKV, r}^\top
\]

\textbf{Advantages:}

\begin{enumerate}
\item \textbf{Shared basis:} $V_{QKV, r}$ is a single input basis for Q, K, V.
\item \textbf{Smaller parameter count:} $r(d + 2d_{kv}) + r \cdot d = 2rd + 2r d_{kv}$ vs
  independent $6rd$.
\item \textbf{Correlation exploitation:} SVD automatically captures the shared
  structure.

\end{enumerate}
We apply separate SVD to the output projection $W_O \in \mathbb{R}^{d \times d_h H}$:

\[
W_O \approx U_{O, r} \Sigma_{O, r}^{1/2} \cdot \Sigma_{O, r}^{1/2} V_{O, r}^\top
\]

\subsection{Contributions}

\begin{itemize}
\item \textbf{Joint QKV SVD.} We propose stacking Q, K, V and factorizing them with a
  single SVD, sharing the input basis across all three projections.
\item \textbf{Multi-architecture support.} The implementation handles flat (DistilBERT),
  nested (RoBERTa), Conv1D (GPT-2), and GQA+RoPE (Qwen) models.
\item \textbf{Correct nested-encoder detection.} We check for \texttt{out\_lin} / \texttt{c\_proj} /
  \texttt{o\_proj} \textbf{before} declaring a model nested, fixing a subtle double-
  projection bug that caused 50.92\% collapse in earlier implementations.
\item \textbf{Safe failure handling.} Weight extraction failures trigger warnings, not
  exceptions, allowing graceful degradation.
\item \textbf{LoRA recovery on both branches.} Both the QKV and output branches have
  independent LoRA adapters.
\item \textbf{Honest parameter accounting.} We report exact parameter reductions.

\end{itemize}
\subsection{Paper Organization}

Section 2 reviews related work. Section 3 provides background on SVD, LoRA, 
and attention. Section 4 presents Compressed MHA. Section 5 describes 
integration. Section 6 details the setup. Section 7 reports results. Section 8 
analyzes trade-offs. Section 9 discusses practical guidance. Section 10 lists 
limitations. Section 11 concludes.



\subsection{Notation}

\begin{lstlisting}

---
\end{lstlisting}